\documentclass{article}
\usepackage{arxiv}

\usepackage[T1]{fontenc}
\usepackage[utf8]{inputenc}
\usepackage{hyperref}
\usepackage{url}
\usepackage{booktabs}
\usepackage{amsfonts}
\usepackage{nicefrac}
\usepackage{microtype}
\usepackage{graphicx}
\usepackage{lipsum}


\title{Radar Posture Classification: A Comparative Study of SOTA Architectures}
\author{Eloan Tourtelier, Othmane Ouali, Lucas Cheylan}
\date{2026}

\begin{document}

\maketitle

\begin{abstract}
We study posture classification from radar point clouds and compare several model
architectures for distinguishing between sitting and standing positions. We describe our radar point representation, training protocol, and evaluation setting, and we benchmark multiple baselines under the same experimental conditions. Our results highlight which architectural choices are most effective for sparse radar point data and provide practical guidance for future radar-based human state recognition.
\end{abstract}

\tableofcontents

\section{Introduction}



\section{Architectures}
\subsection{PointNet}

The PointNet architetcure is a fundamentalm Deep learning architecture applied on radar points, ie unoreder sets of points with their three spatial postions. It was frist made to classify a cloud fo point, and a forking branch of the model can also segment parts of the cloud point to cut it into recognized portions.

The architecture is quite straighforward. It has to be point-order invariants as a cloud point doesnt have point ordering. The netword uses a specific block called a TNet, whihc from a hidden layer computes a matrix to apply to that hidden layer, hence making the point representaiton invariant to roations, sclaing or shearing (whihc is was we want for a point distribution). Those are the main two fetaures of the PointNet. It was origianlly trained for a dataset of 

\begin{figure}[h]
    \centering
    \includegraphics[width=0.6\textwidth]{pointnetarchi.png}
    \caption{The architecture of a PointNet. The intresting layer is the TNet.}
    \label{fig:radar-point-cloud}
\end{figure}


\subsection{Moments}
% Prior work on point clouds, engineered features, sparse radar representations.

\subsection{XGBoost}
% PointNet, PointNet++, PointMLP, transformer-based models, etc.



\section{Data}
\subsection{RadHAR dataset}
% Briefly explain how raw radar is converted into point clouds.

\subsection{Homemade dataset}
% Centering, normalization, padding/sampling to N points, temporal aggregation.

\subsection{Data processing and augmentation}

\section{Results}
\subsection{Comparisons}
Here we compare evrything bout our different approchas, F1 and accuracy scores, amounf of parameters, Inference time etc.

\section{Ablation Studies}

\section{Conclusion}
% Restate findings and next steps.

\nocite{*}
\bibliographystyle{plain}
\bibliography{ref}


\section{Appendix}

\end{document}